\begin{abstract}
Electrocardiogram (ECG) denoising plays a vital role in mobile health and wearable monitoring, where real-world recordings are heavily corrupted by baseline wander (BW), muscle artifacts (MA), and electrode motion (EM). While conventional deep denoising networks achieve substantial noise suppression, they frequently cause diagnostic distortion by blunting high-frequency QRS complexes and attenuating weak P/T waves. Furthermore, denoising models are rarely validated on downstream clinical diagnosis under standardized inter-patient protocols. In this paper, we propose \HSL, an ultra-lightweight 75{,}133-parameter one-dimensional wavelet residual network designed for morphology-preserving ECG reconstruction and automated arrhythmia detection. \HSL{} integrates discrete Haar wavelet decomposition directly into a symmetric U-Net backbone with standard residual blocks with squeeze-and-excitation channel attention (\texttt{ResSEBlock1D}), and high-frequency detail subband residual filtering. To prevent R-peak clipping and baseline distortion, the network is trained with a time-domain mixed loss objective combining mean squared error and absolute error ($\alpha=0.8$). For downstream arrhythmia classification, we introduce \RRN, which couples a 1D ResNet morphology encoder (with SiLU activations) with an independent 6-dimensional RR-interval contextual pathway and interchangeable classical MLP or 6-qubit parameterized variational quantum circuit (VQC) classification heads. Evaluated on the MIT-BIH Arrhythmia Database paired with NSTDB physiological noise across 7 noise combinations and 10 SNR levels ($-10$ to $+10$\,dB), \HSL{} achieves an average SNR improvement ($\Delta$SNR) of \best{10.01\,dB} and a percentage root-mean-square difference (PRD) of \best{33.88\%}, outperforming recent SOTA DeepFilter (9.64\,dB / 35.63\%) and four baseline architectures while requiring only 35.20\,ms inference time on edge CPU. On the clinical gold-standard de Chazal inter-patient DS1/DS2 test split comprising 49{,}298 independent beats from 22 unseen subjects, the downstream classifier achieves an Accuracy of \best{96.43\%}, a Macro-F1 score of \best{0.8407}, and an S-class positive predictivity of \best{72.97\%} (+23.27\% over prior benchmarks) (Classical ensemble) and \best{96.04\%} / \best{0.8176} (Hybrid quantum ensemble). Ablation studies demonstrate that channel-wise attention (SE1D) preserves high-frequency QRS morphology substantially better than spatial attention (CBAM1D), providing an efficient and reliable foundation for wearable cardiac monitoring.
\end{abstract}

\begin{IEEEkeywords}
Electrocardiogram (ECG) denoising, discrete wavelet transform (DWT), squeeze-and-excitation, heartbeat classification, variational quantum circuit (VQC), inter-patient evaluation, wearable edge computing.
\end{IEEEkeywords}
